\writeSectionFrame{\presentationTitle}{Einführung}
\note{
    Da Entwurfsmuster sprachunabhängig in UML-Diagrammen beschrieben werden, wiederholen wir noch einmal die wichtigsten UML-Diagramme.
}

\begin{frame}{Was ist UML?}
    \begin{itemize}
        \item UML = Unified Modeling Language
        \item Grafische Sprache für Software-Modelle
        \item Standardisiert und weit verbreitet
    \end{itemize}
    \note{
        UML ist eine standardisierte Modellierungssprache, die zur Visualisierung, Spezifikation, Konstruktion und Dokumentation von Software-Systemen verwendet wird. Sie spielt eine zentrale Rolle in der objektorientierten Programmierung, indem sie die Beziehungen und Interaktionen zwischen Klassen und Objekten darstellt.
    }
\end{frame}

\section{Warum UML verwenden?}
\begin{frame}{Warum UML verwenden?}
    \begin{itemize}
        \item Standardisierte Notation
        \item Verbessert Kommunikation
        \item Unterstützt Planung und Abstraktion
    \end{itemize}
    \note{
        UML bietet eine einheitliche, international akzeptierte Methode zur Beschreibung von Softwaresystemen, die es erleichtert, mit verschiedenen Stakeholdern zu kommunizieren. Es hilft Entwicklern, Architekten und Analysten, auf einer gemeinsamen Basis zu arbeiten. Durch Abstraktion können komplexe Systeme besser geplant und strukturiert werden.
    }
\end{frame}

\section{Geschichte der UML}
\begin{frame}{Geschichte der UML}
    \begin{itemize}
        \item Entwickelt in den 1990er Jahren
        \item Gründer: Booch, Jacobson, Rumbaugh
        \item Standard der Object Management Group (OMG)
    \end{itemize}
    \note{
        UML wurde in den frühen 1990er Jahren von Grady Booch, Ivar Jacobson und James Rumbaugh entwickelt. Ihr Ziel war es, die besten Konzepte aus verschiedenen objektorientierten Modellierungstechniken zu vereinen und eine einzige, standardisierte Sprache zu schaffen. UML wird heute von der OMG (Object Management Group) gepflegt und weiterentwickelt.
    }
\end{frame}

\section{Überblick über UML-Elemente}
\begin{frame}{Überblick über UML-Elemente}
    \begin{itemize}
        \item Zwei Hauptkategorien: Struktur- und Verhaltensdiagramme
        \item Verschiedene Diagramme für verschiedene Zwecke
    \end{itemize}
    \note{
        UML besteht aus verschiedenen Diagrammtypen, die unterschiedliche Aspekte eines Systems modellieren. Es gibt zwei Hauptkategorien: Strukturdiagramme, die statische Elemente wie Klassen, Komponenten und Objekte beschreiben, und Verhaltensdiagramme, die dynamische Aspekte wie Interaktionen, Zustände und Aktivitäten darstellen.
    }
\end{frame}

\section{Wer verwendet UML?}
\begin{frame}{Wer verwendet UML?}
    \begin{itemize}
        \item Entwickler, Architekten, Analysten
        \item Projektmanager und Qualitätssicherer
    \end{itemize}
    \note{
        UML wird von einer Vielzahl von Personen verwendet, die an der Softwareentwicklung beteiligt sind. Entwickler nutzen es zur Modellierung von Klassen und Interaktionen. Architekten verwenden UML zur Planung von Systemstrukturen, während Geschäftsanalysten es zur Modellierung von Prozessen und Anforderungen verwenden. Projektmanager und Qualitätssicherer nutzen UML zur Überprüfung des Fortschritts und der Richtigkeit von Modellen.
    }
\end{frame}

\section{Vorteile von UML}
\begin{frame}{Vorteile von UML}
    \begin{itemize}
        \item Klarheit und Präzision
        \item Skalierbar für kleine und große Systeme
        \item Standardisierung und Flexibilität
    \end{itemize}
    \note{
        UML bietet eine klare und präzise Möglichkeit, komplexe Systeme zu beschreiben. Es ist flexibel und kann sowohl für kleine als auch große Systeme verwendet werden. Die Standardisierung von UML erleichtert die Zusammenarbeit über Teams und Unternehmen hinweg. Es gibt viele Werkzeuge, die die Erstellung und Bearbeitung von UML-Diagrammen unterstützen, was es einfacher macht, Modelle zu erstellen und zu aktualisieren.
    }
\end{frame}

\begin{frame}{UML-Diagrammtypen - Grobe Einteilung}
    \begin{itemize}
        \item Strukturdiagramme (statisch)
        \item Verhaltensdiagramme (dynamisch)
    \end{itemize}
    \note{
        UML-Diagramme sind in zwei Hauptkategorien unterteilt: Strukturdiagramme und Verhaltensdiagramme. 
        \newline
        \textbf{Strukturdiagramme} zeigen die statischen Aspekte eines Systems, wie die Klassenstruktur, Objekte, Komponenten oder die Verteilung von Software auf Hardware.
        \newline
        \textbf{Verhaltensdiagramme} modellieren dynamische Aspekte, wie Use Cases, Aktivitäten, Zustände und Interaktionen zwischen Objekten.
    }
\end{frame}

\begin{frame}{Strukturdiagramme}
    \begin{itemize}
            \item Klassendiagramm
            \item Objektdiagramm
            \item Komponentendiagramm
            \item Kompositionsstrukturdiagramm
            \item Verteilungsdiagramm
            \item Paketdiagramm
    \end{itemize}
\end{frame}

\begin{frame}{Verhaltensdiagramme}
    \begin{itemize}
        \item Anwendungsfalldiagramm
        \item Aktivitätsdiagramm
        \item Zustandsdiagramm
        \item Sequenzdiagramm
        \item Kommunikationsdiagramm
        \item Interaktionsübersichtsdiagramm
        \item Zeitverlaufsdiagramm
    \end{itemize}
\end{frame}